\section{The rotated-pentagon profile}
\label{sec:rotated-regular-polygons}
\providecommand{\pentagonassetpath}{figures/}

The Haim--Kislev--Ostrover counterexample is a Lagrangian product of
two regular pentagons. We determine how its capacity changes when one
pentagon rotates relative to the other. The volume stays constant, so
this also determines which relative positions maximize the systolic ratio.

Let $P_5\subset\mathbb R^2$ be a regular pentagon centred at the origin
with circumradius one. Write $R(\theta)$ for counterclockwise rotation
through angle $\theta$, and put
\[
 K_\theta=P_5\times_L R(\theta)P_5
 \subset\mathbb R_q^2\oplus\mathbb R_p^2.
\]
The Lagrangian product here is the Cartesian product in the two coordinate
planes, with symplectic form
$\omega_0((q,p),(q',p'))=q\cdot p'-p\cdot q'$.
Simultaneous rotation of both planes is symplectic, so the choice of
initial orientation for $P_5$ does not affect the capacity. We use
\[
 \operatorname{sys}(K)=
 \frac{c_{\mathrm{EHZ}}(K)^2}{2\operatorname{vol}_4(K)}.
\]

The sampled ratios in Figure~\ref{fig:rotated-pentagons-curve} rise from
the aligned position to a maximum at $\theta=\pi/10$, then fall
symmetrically. They suggested a profile proportional to $\sec^2\theta$
on $[0,\pi/10]$. We prove this formula by comparing the quadratic
programs at different angles. Their feasible weights remain the same
under rotation, and each fixed choice has a particularly simple angular
dependence. The known capacity at the Haim--Kislev--Ostrover position
will then bound every such choice throughout the interval.

\begin{figure}[htbp]
 \centering
 \includegraphics[width=0.82\linewidth]{\pentagonassetpath rotation-profile.png}
 \caption{Systolic ratios for the rotating pentagon product. The blue
 markers are numerical evaluations of the finite quadratic programs;
 the black curve is the formula proved in this chapter.
 The horizontal line marks $\operatorname{sys}=1$.}
 \label{fig:rotated-pentagons-curve}
\end{figure}

To state the profile for all angles, define
\[
 d(\theta)=\min_{k\in\mathbb Z}\left|\theta-\frac{k\pi}{5}\right|.
\]
This is the distance to the nearest multiple of $\pi/5$, and takes
values in $[0,\pi/10]$.

\begin{theorem}[Rotated-pentagon profile]
\label{thm:rotated-pentagon-product-formula}
For every real $\theta$,
\begin{align}
 c_{\mathrm{EHZ}}(K_\theta)
 &=\bigl(1+\cos(\pi/5)\bigr)^2\sec d(\theta),
 \label{eq:pentagon-capacity-profile}\\
 \operatorname{sys}(K_\theta)
 &=\frac{5+2\sqrt5}{10}\sec^2 d(\theta).
 \label{eq:pentagon-systolic-profile}
\end{align}
The maximum systolic ratio is $(3+\sqrt5)/5>1$, attained exactly at
$\theta=\pi/10+k\pi/5$, $k\in\mathbb Z$.
\end{theorem}

\subsection{Symmetries and the Haim--Kislev--Ostrover position}

We first reduce the angles to $[-\pi/10,\pi/10]$ and identify the
capacity at its endpoints. The pentagon's rotational symmetry gives
$K_{\theta+2\pi/5}=K_\theta$. If $S$ is reflection in a symmetry axis
of $P_5$, the simultaneous reflection $(q,p)\mapsto(Sq,Sp)$ is
symplectic and sends $K_\theta$ to $K_{-\theta}$. Thus capacity is even
as a function of $\theta$.

Exchanging the two factors gives a further symmetry. The symplectic map
$(q,p)\mapsto(p,-q)$ sends $K_\theta$ to
$R(\theta)P_5\times_L(-P_5)$. Rotate both factors through $-\theta$
and use $-P_5=R(\pi/5)P_5$ to obtain $K_{\pi/5-\theta}$.
Together with evenness, this proves that capacity has period $\pi/5$.
Every angle can therefore be reduced to $[-\pi/10,\pi/10]$ without
changing capacity; evenness will give the final expression in terms of
$d(\theta)$.

Haim--Kislev and Ostrover computed the capacity of the product with
relative angle $-\pi/2$. Their convention for the symplectic form is
$\sum_i dp_i\wedge dq_i=-\omega_0$. The map $(q,p)\mapsto(q,-p)$
changes this convention to ours and sends their body to $K_{\pi/2}$.
Since $\pi/2=\pi/10+2\pi/5$, their result
\cite[Proposition~1.4]{HaimKislevOstrover2024}, together with the reflection
symmetry, gives
\begin{equation}
 c_{\mathrm{EHZ}}(K_{-\pi/10})
 =c_{\mathrm{EHZ}}(K_{\pi/10})
 =2\cos(\pi/10)\bigl(1+\cos(\pi/5)\bigr).
 \label{eq:pentagon-known-endpoints}
\end{equation}

\subsection{Feasible weights under rotation}

The quadratic capacity formula lets us compare different angles using
the same facet order and weights. We spell out its variables for this
family. Set
\[
 h=\cos(\pi/5),\qquad
 n_i=\left(\cos\left(\frac\pi2+\frac{2\pi i}{5}\right),
           \sin\left(\frac\pi2+\frac{2\pi i}{5}\right)\right),
 \quad 0\le i<5.
\]
Choose the orientation $P_5=\{q:n_i\cdot q\le h,\ 0\le i<5\}$.
The ten facets of $K_\theta$ then have normalized normals
\[
 a_i(\theta)=h^{-1}(n_i,0),\qquad
 a_{i+5}(\theta)=h^{-1}(0,R(\theta)n_i),\qquad 0\le i<5,
\]
so that their inequalities are $a_i(\theta)\cdot x\le1$.

A facet word $\sigma=(\sigma_1,\ldots,\sigma_m)$ is an ordered list
of distinct indices from $\{0,\ldots,9\}$. Its feasible weights satisfy
\begin{equation}
 \beta_k\ge0,\qquad \sum_{k=1}^m\beta_k=1,\qquad
 \sum_{k=1}^m\beta_k a_{\sigma_k}(\theta)=0.
 \label{eq:pentagon-feasibility}
\end{equation}
The last equation is closure. For these weights, define
\begin{equation}
 Q_{\sigma,\beta}(\theta)=
 \sum_{1\le j<k\le m}\beta_j\beta_k
 \omega_0\bigl(a_{\sigma_j}(\theta),a_{\sigma_k}(\theta)\bigr).
 \label{eq:pentagon-q-definition}
\end{equation}
This uses the traversal order of the quadratic-program chapter.
The Haim--Kislev formula~\cite{HK2017} states that
\begin{equation}
 c_{\mathrm{EHZ}}(K_\theta)=\frac1{2Q_{\max}(\theta)},
 \qquad
 Q_{\max}(\theta)=\max_{\sigma,\beta}Q_{\sigma,\beta}(\theta)>0,
 \label{eq:pentagon-hk}
\end{equation}
where the maximum is over all words and their feasible weights.
There are finitely many words, and each feasible set is a closed subset
of a simplex, so the maximum exists. Zero weights allow facets to be
omitted.

For this product, closure separates into equations in the two planes:
\[
 \sum_{\sigma_k<5}\beta_k n_{\sigma_k}=0,
 \qquad
 R(\theta)\sum_{\sigma_k\ge5}\beta_k n_{\sigma_k-5}=0.
\]
Removing the invertible rotation from the second equation shows that
feasibility is independent of $\theta$. Thus a word and weights chosen
at one angle remain feasible at every other angle.

We first exhibit a choice that gives the proposed capacity as an upper
bound. Take
\[
 \sigma=(3,8,1,0,5,6),\qquad
 \beta=\frac1{4(1+h)}(2h,2h,1,1,1,1).
\]
These weights are positive and sum to one. The normal $n_3$ points
opposite the bisector of $n_0$ and $n_1$, and hence
\[
 n_0+n_1=-2h n_3.
\]
This identity gives closure in both factors. To evaluate $Q$, group the
consecutive entries $(1,0)$ and $(5,6)$ of the word. Normals from the
same factor have zero symplectic pairing, so bilinearity lets us replace
each group by its weighted sum. Writing
\[
 u=\frac{n_3}{2(1+h)},\qquad v=R(\theta)u,
\]
the four weighted sums, in their order, are
\[
 (u,0),\quad(0,v),\quad(-u,0),\quad(0,-v).
\]
Their ordered symplectic pairings sum to $2u\cdot v$. Therefore
\begin{equation}
 Q_{\sigma,\beta}(\theta)
 =\frac{\cos\theta}{2(1+h)^2}.
 \label{eq:pentagon-attaining-word}
\end{equation}
For $|\theta|\le\pi/10$ this is positive, and the capacity formula yields
\begin{equation}
 c_{\mathrm{EHZ}}(K_\theta)\le(1+h)^2\sec\theta.
 \label{eq:pentagon-upper-bound}
\end{equation}
To prove equality, we must show that no other feasible choice gives a
larger quadratic value.

\subsection{Bounding all choices by their endpoint values}

Fix any facet word $\sigma$ and feasible weights $\beta$.
Same-factor pairings in~\eqref{eq:pentagon-q-definition} vanish, and
each mixed pairing is, up to sign, $h^{-2}n_i\cdot R(\theta)n_j$.
There are consequently constants $A_{\sigma,\beta}$ and
$B_{\sigma,\beta}$ such that
\[
 Q_{\sigma,\beta}(\theta)=
 A_{\sigma,\beta}\cos\theta+B_{\sigma,\beta}\sin\theta.
\]
This angular dependence determines the value inside our interval from
the two endpoint values. The sine addition formulas give
\begin{align}
 Q_{\sigma,\beta}(\theta)
 &=\frac{\sin(\pi/10-\theta)}{\sin(\pi/5)}
       Q_{\sigma,\beta}(-\pi/10)\notag\\
 &\quad+\frac{\sin(\pi/10+\theta)}{\sin(\pi/5)}
       Q_{\sigma,\beta}(\pi/10).
 \label{eq:pentagon-interpolation}
\end{align}
For $-\pi/10\le\theta\le\pi/10$, both coefficients are nonnegative
and their sum is $\cos\theta/\cos(\pi/10)$.

The fixed choice remains feasible at both endpoints, so its quadratic
value there is at most the corresponding maximum. By
\eqref{eq:pentagon-known-endpoints} and~\eqref{eq:pentagon-hk},
\[
 Q_{\sigma,\beta}(\pm\pi/10)
 \le Q_{\max}(\pm\pi/10)
 =\frac1{4\cos(\pi/10)(1+h)}.
\]
We may substitute these upper bounds into
\eqref{eq:pentagon-interpolation} because its coefficients are
nonnegative. This gives
\[
 Q_{\sigma,\beta}(\theta)
 \le\frac{\cos\theta}{4\cos^2(\pi/10)(1+h)}
 =\frac{\cos\theta}{2(1+h)^2},
\]
using $2\cos^2(\pi/10)=1+h$. The bound holds for every feasible
choice, and the choice in~\eqref{eq:pentagon-attaining-word} attains it.
We have proved
\[
 Q_{\max}(\theta)=\frac{\cos\theta}{2(1+h)^2},\qquad
 c_{\mathrm{EHZ}}(K_\theta)=(1+h)^2\sec\theta,
 \qquad |\theta|\le\pi/10.
\]
Evenness and $\pi/5$-periodicity give
\eqref{eq:pentagon-capacity-profile} for all real angles.

Finally, the five triangles joining the centre of $P_5$ to its edges
give
\[
 \operatorname{area}(P_5)=\frac52\sin(2\pi/5),\qquad
 \operatorname{vol}_4(K_\theta)=\left(\frac52\sin(2\pi/5)\right)^2.
\]
Substitution into the definition of the systolic ratio yields
\[
 \operatorname{sys}(K_\theta)
 =\frac{(1+\cos(\pi/5))^4}
        {2\left(\frac52\sin(2\pi/5)\right)^2}\sec^2 d(\theta)
 =\frac{5+2\sqrt5}{10}\sec^2 d(\theta).
\]
Since $\sec^2 t$ is strictly increasing on $[0,\pi/10]$, the maximum
occurs precisely when $d(\theta)=\pi/10$. At those angles its value is
$(3+\sqrt5)/5$. This completes the proof of
Theorem~\ref{thm:rotated-pentagon-product-formula}.

% Proof provenance: the endpoint interpolation argument is supplied in
% docs/review-evidence/pro-review-2026-09-15/thesis_review/
% 03_mathematical_audit.md, section 2. The explicit word and numerical
% rotation profile predate that review in thesis/09-rotated-regular-polygons*.
